\section{Introduction and main results}\label{art:introduction}

Can ordinary BFGS converge locally on a convex objective whose gradients jump?
Regularity of the formulas defining the objective does not by itself settle
this question. We prove local convergence when a finite polyhedral arrangement
controls the nearby smooth pieces, and construct local failure when four smooth
convex formulas encode an accumulating sequence of bad switching configurations.
The distinction is between the geometry of switching and the regularity within
individual branches. It persists when all gradients queried by the algorithm
are ordinary derivatives.

Throughout, BFGS means the unmodified full inverse update. An infinite
\emph{admissible run} starts with any positive definite matrix, uses fixed
$0<c_1<c_2<1$, satisfies Armijo decrease and weak Wolfe curvature, and accepts
only differentiable points with nonzero gradients; the equations are given in
Definition~\ref{art:admissible}. This notion describes successful successive
line searches. It does not guarantee that every search ray contains a
differentiable acceptable point. The negative results instead verify every
decision of one specified unit-first doubling/bisection implementation.

The following theorem gives the main positive contrast. The local extensions
are defined on open sets, so smoothness at cell boundaries has a precise meaning.
\begin{theorem}[Local convergence on finite polyhedral cells]\label{art:polyhedral}
Let $n\ge1$ and let $F:\R^n\to\R$ be finite and convex, with a unique minimizer $x_*$.
Suppose an open neighborhood $U$ of $x_*$ is covered by finitely many closed
full-dimensional polyhedra $P_i$. For each $i$, suppose there is an open set
$V_i\supset U\cap P_i$ and a function $f_i\in C^{1,1}_{\rm loc}(V_i)$ such that
\[
 F(x)=f_i(x)\qquad(x\in U\cap P_i).
\]
Then there exists $\varepsilon>0$ such that every infinite full-BFGS run with
arbitrary $H_0\succ0$, differentiable accepted iterates, fixed
Armijo--weak-Wolfe constants $0<c_1<c_2<1$, and
$F(x_0)\le F(x_*)+\varepsilon$ satisfies
\[
 x_k\in U\quad(k\ge0),\qquad F(x_k)\to F(x_*),\qquad x_k\to x_*.
\]
The basin is chosen before the run and its initial matrix. Only the full
objective must be convex; the extensions need not be convex outside their cells.
\end{theorem}


Finite maxima of smooth functions need not satisfy this finite-cell hypothesis.
The next result fixes the objective and implementation before choosing a
sequence of failing initial points. Its proof constructs all trial points,
not only an accepted orbit.
\begin{theorem}[Local failure with globally smooth convex branches]
\label{art:local-failure}
There are four globally convex functions $f_i\in C^{2,1}(\R^2)$ such that
\[
 0\preceq\nabla^2f_i(x)\preceq64I\qquad(i=0,1,2,3; x\in\R^2)
\]
and $F=\max_i f_i$ has compact sublevel sets and a unique nonsmooth
minimizer at zero. For this one objective there are fixed starting points
$x_0^{(\ell)}\to0$, all with $H_0=I$, such that, for every fixed pair
\[
 0<c_1\le\tfrac13,\qquad c_1<c_2<1,
\]
the deterministic unit-first bracketing implementation of Section~\ref{art:failure} produces infinite full-BFGS runs satisfying
\[
 x_k^{(\ell)}\to \bar x_\ell\ne0,\qquad
 F(x_k^{(\ell)})\downarrow F(\bar x_\ell)>F(0)=0,\qquad
 0\notin\partial F(\bar x_\ell).
\]
Every queried point is differentiable. Each line search takes at most five
trials, and its accepted lengths repeat $1,1/4,1/16$.
The objective, initial points, initial matrix and complete decision tree
are the same throughout the displayed parameter range. The supremum of
the norms of all accepted and rejected query points in the $\ell$th run
tends to zero as $\ell\to\infty$.
\end{theorem}


Theorem~\ref{art:local-failure} rules out a local guarantee based only on a
finite number of very smooth convex branch formulas. It does not say that every
nearby starting point fails. The obstruction is compatible with compact
sublevels and a unique minimizer. An independent construction changes the
balance between smoothness and lower curvature.
\begin{theorem}[Independent local failure with positive lower curvature]
\label{art:local-failure-lower}
There is a finite strongly convex function $F:\R^2\to\R$ with compact
sublevels and a unique nonsmooth minimizer at zero, with fixed starting
points $x_0^{(\ell)}\to0$, all with $H_0=I$, such that for every fixed pair
$0<c_1\le1/3$, $c_1<c_2<1$ the deterministic unit-first bracketing
implementation of Section~\ref{art:failure} produces infinite
full-BFGS runs, the same for every such pair. They converge to nonstationary points $\bar x_\ell\ne0$,
with $\bar x_\ell\to0$ and
$F(x_k^{(\ell)})\downarrow F(\bar x_\ell)>F(0)$. Every accepted and rejected query
is differentiable, the accepted lengths repeat $1,1/4,1/16$, and the
supremum of all query norms in the $\ell$th run tends to zero.

The objective can be a maximum of four globally $C^{1,1}$, uniformly
strongly convex branches with globally Lipschitz gradients. There are
fixed $0<\underline\mu\le\overline L<\infty$ such that
\[
 \underline\mu I\preceq\nabla^2F(x)\preceq\overline L I
 \quad\text{at every twice-differentiable point}.
\]
The first three branches in this construction are not twice
differentiable at zero.
\end{theorem}


In this last statement, a classical Hessian means the derivative of the gradient
map, with that map defined on a neighborhood of the point. An isolated second
Peano expansion is not the convention used. Table~\ref{art:regularity-map}
keeps the incomparable regularity conclusions separate. In particular, neither
construction gives local failure for a finite maximum of uniformly strongly
convex $C^2$ branches through the optimizer. The common quadratic example used
inside the second proof fails globally near a wrong corner; it is the subsequent
embedding that gives failure near its own minimizer.

\begin{table}[t]
\centering
\begin{tabular}{p{.27\textwidth}p{.31\textwidth}p{.31\textwidth}}
\toprule
Property & Theorem~\ref{art:local-failure} & Theorem~\ref{art:local-failure-lower}, finite-branch version\\
\midrule
Number of branches & Four & Four\\
Branch regularity & Global $C^{2,1}$ & Global $C^{1,1}$\\
Positive lower curvature & Not imposed & Uniformly strongly convex\\
At the accumulation point & Branch Hessians exist & Three branch Hessians fail to exist\\
Initialization and locality & Fixed $H_0=I$; all queries approach the optimizer & Fixed $H_0=I$; all queries approach the optimizer\\
\bottomrule
\end{tabular}
\caption{Two local counterexamples along different regularity axes. The last
row concerns a sequence of runs for one fixed objective in each column.}
\label{art:regularity-map}
\end{table}

The positive proof starts from a comparison energy. Its negative term measures
progress still visible to a supporting gradient; its positive term measures a
secant correction. We establish a hierarchy for the actual principal minors
of full-BFGS matrices and derive divergence of the reciprocal arithmetic mean
of $g_k^TH_kg_k$. This provides a budget that excludes a persistent objective
gap when the objective supplies a suitable comparison field. The direction
toward a fixed minimizer gives the finite-cell result. The criterion is sufficient, and is
not claimed to characterize all convergent BFGS runs.

\subsection{Related work and the role of the present results}
Classical BFGS analysis combines convexity with matrix potentials and
line-search inequalities. Powell's convergence analysis, the role of the
quasi-Newton matrices studied by Byrd--Nocedal, and Nocedal's survey provide
the relevant smooth background~\cite{FND_P76,FND_BN89,FND_N92}.
The Gram and Schur identities used below are standard matrix algebra.
Our use of adjacent principal-minor degrees supplies a bridge from the full
matrix update to a reciprocal budget in the bounded-gradient, vanishing-decrease
regime. We do not attribute the underlying determinant, trace, or AM--GM tools
to this work.

For unmodified nonsmooth BFGS, Guo--Lewis proves convergence on the Euclidean
norm in every finite dimension~\cite[Corollary~4.2]{GL18}; this is a known
special case of the centered homogeneous setting below. Lewis--Overton studies
the practical method and the line search used in our examples~\cite{LO}.
The finite two-piece affine escape theorem of
\preprintcite[Theorem IV.1]{XW} is also prior work, with an unbounded-descent
conclusion rather than the compact-level conclusion considered here.
Gebken's criticality result assumes an iterate limit and specified matrix
behavior~\cite{BG}; those are not premises in our convergence theorems.

Earlier negative examples have different combinations of objective geometry
and algorithmic hypotheses. Smooth nonconvex BFGS failures~\cite{NEG_M04,NEG_M07},
the examples of \preprintcite{NEG_LO08}, and other nonconvex quasi-Newton
constructions~\cite{NEG_LZ15} are relevant antecedents, but do not establish
the two convex local statements above. Limited-memory failure
results~\cite{NEG_AO21} concern a different update. The recent smooth
nonconvex planar example of \preprintcite{NEG_DIAO26} likewise has different
objective assumptions. We compare these statements by their hypotheses;
no claim of failure for every weak-Wolfe implementation is made.

The geometric interpretation of surviving smooth directions is related to
$VU$ theory, partial smoothness, and stratification
\cite{geom_Lewis2002,LL,BDLS}. Convex-gradient cocoercivity is an established
ingredient~\cite{WW}. The companion develops affine-critical and fixed-normal
extensions using those antecedents. Such connections identify standard tools
and related structures, without supplying a necessary-and-sufficient dynamical
classification or establishing priority for a complete theorem.

Sections~\ref{art:geometry}--\ref{art:positive} develop the geometric
motivation and the positive proof. Section~\ref{art:beyond} identifies the
separate companion results and their limits. Section~\ref{art:failure} proves
Theorem~\ref{art:local-failure}, with its finite certificate in
Appendix~\ref{art:seed_certificate}; Appendix~\ref{art:lower_curvature}
proves Theorem~\ref{art:local-failure-lower}. All essential arguments for the
article's theorems are contained here, apart from the explicitly invoked
standard background theorems.
